\chapter{Real-time data handling}\label{ch:realtime}

The Live page and the future alert service both consume a continuous stream. This chapter describes how that stream is
received, buffered and kept alive, and how the other path, the download of day files, works.

\section{SeedLink}
SeedLink is a protocol for distributing seismic data over TCP. A client connects to the server (on a Shake, port~18000),
asks which stations and streams it offers (\code{INFO STREAMS}), selects the channels it wants, and then simply
receives miniSEED records as they are completed by the digitiser. The application uses ObsPy's
\code{EasySeedLinkClient} for the protocol.

\begin{rs3dapp}
\file{core/seedlink.py} first \emph{probes} the Shake (connect, parse the XML answer of \code{INFO STREAMS}, disconnect)
to learn the network, station, location and channel names, which is what the Status page and \ui{Find Shake} use. The
streaming class \code{SeedLinkSource} then selects \code{EHZ}, \code{EHN} and \code{EHE}. ObsPy reconnects by itself
inside its collect loop, but its defaults (120~s idle timeout, 30~s delay) are far too slow for a live display, so the
application lowers them to a 10~s network timeout and a 2~s delay. One workaround is documented in the source: the
client is created with \code{autoconnect=False} and connected manually, because ObsPy~1.5.1 crashes when a timeout is
set after an automatic connection.
\end{rs3dapp}

\section{The ring buffer}
Incoming records are appended to a \emph{ring buffer} that holds the most recent 10~minutes per channel. Records arrive
in short pieces (typically a few seconds each); the buffer stores each channel as a list of \emph{contiguous chunks}:
a new piece is appended to the last chunk when it continues it exactly (same sampling rate, start within half a
sample interval of the previous end) and otherwise starts a new chunk, so every chunk is gap-free by construction.
Chunks that end more than 600~s before the newest data are discarded. All access is guarded by a lock, because
the receiving thread writes while the display thread reads.

To draw, the display asks for a \emph{snapshot}: the samples of one channel in the window $[t_e-W,\,t_e]$, where $t_e$
is the newest time and $W$ the chosen window (30~s to 10~min). The snapshot is assembled from the chunks, with a
\code{NaN} separator between chunks so that gaps stay visible. If a band-pass is requested it is applied to each chunk on
its own (Chapter~\ref{ch:filter}); then, if more than 1500 points remain, min/max decimation (Chapter~\ref{ch:vis})
reduces them. The page redraws every 0.25~s, and \ui{Pause} freezes the picture while the buffer continues to fill.

\section{Keeping the link alive}
Networks fail: a cable is unplugged, a Shake restarts. The streaming thread therefore runs a small state machine
(\cref{fig:reconnect}). It starts in \emph{connecting}; when the first record arrives it is \emph{live}; when the
connection ends or fails it waits and tries again in \emph{reconnecting}. The wait grows with each consecutive failure along the
sequence
\begin{equation}
  \text{back-off} = 2,\ 4,\ 8,\ 16,\ 30\ \text{seconds}\quad(\text{then 30 s}),
\end{equation}
so a Shake that is off for an hour is not polled every two seconds, and the sequence restarts at 2~s as soon as a connection
has delivered data. A second condition covers the case where the TCP connection survives but data stop: if no record
has arrived for 5~seconds while the state is \emph{live}, the page shows \emph{waiting for data}. When a connection
error is the reason, the reason text (for example \emph{connection refused}) is shown next to \emph{reconnecting}.

\begin{figure}[tbp]
\centering
\begin{tikzpicture}[every node/.style={font=\small\sffamily}]
  \node[box] (stop) at (0,0) {stopped};
  \node[box] (conn) at (4.2,0) {connecting};
  \node[boxblue] (live) at (10.6,0) {live};
  \node[box,fill=rs3damber!20,draw=rs3damber] (rec) at (6.6,-2.3) {reconnecting};
  \draw[flow] (stop) -- node[above]{start} (conn);
  \draw[flow] (conn) -- node[above]{first record} (live);
  \draw[flow] (conn.south) -- ++(0,-0.6) |- node[pos=0.25,left]{error} (rec.west);
  \draw[flow] (live.south) -- ++(0,-0.6) |- node[pos=0.25,right]{link lost} (rec.east);
  \draw[flow] (rec.north) -- node[right,pos=0.5,align=left]{wait 2, 4, 8, 16, 30~s,\\then retry} (6.6,-0.9) -| (conn.south east);
  \draw[flow,dashed] (live.north west) to[out=110,in=70,looseness=7] node[above=1pt]{no data for 5~s: ``waiting''} (live.north east);
\end{tikzpicture}
\caption{States of the live connection. The back-off sequence resets when a connection has delivered data.}\label{fig:reconnect}
\end{figure}

\begin{rs3dapp}
\code{RingBuffer}, \code{LiveStream} and the constants \code{BACKOFF}$=(2,4,8,16,30)$ are in \file{core/live.py}; the
display loop is in \file{app/pages/live.py}. Data are not persisted by the live path: the Live page shows what the buffer
holds; permanent data come from the Shake's own archive through the Data page.
\end{rs3dapp}

\section{Downloading the archive}
The Shake writes its own SDS archive. The application copies day files with SFTP over the same SSH login that the Status
page uses. For each selected channel it lists the remote days, compares them with the days already on the laptop, and fetches
the missing ones \emph{plus the newest remote day}, which is still growing. A file is first written as
\file{<name>.part} and renamed when complete, so a transfer cut by a cable fault leaves no corrupt day file. Reading
is read-only on the Shake; the application never modifies it.

\section{Latency}
The delay between ground motion and the picture has three parts: the digitiser must complete a record before sending
it (a few seconds), the network adds milliseconds on a direct cable, and the page refreshes every 0.25~s. A practical
latency of a few seconds is therefore expected, and is set by the record length, not by the application's processing.
